\documentclass[10pt,a4paper]{amsart}
\usepackage[margin=19mm]{geometry}
\usepackage{amsmath,amssymb,mathtools}
\usepackage[scr=rsfs,cal=euler]{mathalfa}
\usepackage{xfrac}
\usepackage[unicode,hidelinks,bookmarks=false]{hyperref}
\newcommand{\ie}{i.e.\ }
\newcommand{\cf}{cf.\ }
\newcommand{\resp}{respectively\ }
\newcommand{\Subsection}{\subsection}
\providecommand{\nobreakdash}{\nobreak\hbox{-}}
\setlength{\emergencystretch}{2em}
\multlinegap0pt
\usepackage{amssymb}
\usepackage{tikz}
\usepackage{mathtools}
\newtheorem{theorem}{Theorem}
\newtheorem{lemma}{Lemma}
\newcommand{\Cl}{\mathrm{Cl}}
\newcommand{\rk}{\mathrm{rk}}
\title{The existence of infinitely many cubic fields with class group of exact $2$-rank $1$}
\author{Manjul Bhargava, Arul Shankar, Artane Siad, Ashvin Swaminathan}
\date{Source: arXiv:2602.06894v1 (CC BY 4.0)}
\begin{document}
\maketitle

\noindent\emph{Source and licence.} The existence of infinitely many cubic fields with class group of exact $2$-rank $1$. Version arXiv:2602.06894v1 (CC BY 4.0), distributed by arXiv under the Creative Commons Attribution 4.0 International licence, http://creativecommons.org/licenses/by/4.0/, as recorded in the arXiv licence metadata for this exact version and shown on the abstract page https://arxiv.org/abs/2602.06894v1. Full licence notice: \texttt{licenses/arxiv-2602.06894-CC-BY-4.0.txt}.\par\smallskip

\noindent\textit{Editorial scope.} Counted unit: the article's theorem theorem (source0.tex lines 183--185). The article's complete original proof of it is delivered with it (source0.tex lines 245--247, one \texttt{proof} environment of 3 lines, which the article places away from the statement). No mathematics is written, repaired or re-derived: the statement and the proof are the article's own lines, in the article's own notation, and no retained line is substituted or re-flowed.  Retained uncounted as the source-local context the proof needs: lines 229--231; lines 238--240.  Nothing was omitted from any retained span, so the package carries no omission record.  No retained line contains a citation, so the wrapper carries no bibliography and no bibliography entry.  No retained line contains a figure environment, an image-inclusion command or drawing code, so no image is delivered and the package declares no asset dependency.  Editorial formatting only, in addition: an AMS \texttt{amsart} preamble that restates the article's own declarations and macros used by the retained lines, namely the environments \texttt{theorem}, \texttt{lemma} on separate counters, together with the standard packages those lines need.  The retained environments print in this document as Theorem 1, Lemma 1 in order of appearance, whereas the article prints the same items with the numbering of its own full document.  The complete native source of the article as distributed in the arXiv e-print for this exact version is bundled unchanged as \texttt{sources/194136/source0.tex}.
\begin{lemma} \label{lem: finite repetition}
    Up to isomorphism, a cubic field can occur at most $60$ times as the field $\mathbb{Q}[x]/(f(x))$ associated to a maximal polynomial of the form $f(x) = x^3+ax^2+bx+1$ with $a,b \in \mathbb{Z}$. 
\end{lemma}

\begin{lemma}\label{lem: infinitely many have rk_2 = 1}
    There are infinitely many members of the family $+B_{1,1}^2$ that have $\rk_2 \, \Cl = 1$.
\end{lemma}

\begin{theorem} \label{thm: main theorem}
    There are infinitely many cubic fields with class groups of exact $2$-rank $1$.
\end{theorem}

\begin{proof}[Proof of Theorem \ref{thm: main theorem}]
Because infinitely many members of $+B_{1,1}^2$ have class groups of exact $2$-rank $1$ by Lemma \ref{lem: infinitely many have rk_2 = 1} and any cubic field can be repeated as a member of $+B_{1,1}^2$ at most finitely many times by Lemma \ref{lem: finite repetition}, there are infinitely many totally real unit-monogenised cubic fields of exact $2$-rank $1$ and Theorem \ref{thm: main theorem} follows. 
\end{proof}

\end{document}
